\documentclass[a4paper,11pt]{article}
\usepackage[utf8]{inputenc}
\usepackage[T1]{fontenc}
\usepackage[french]{babel}
\usepackage[headheight=15pt, top=2.5cm, bottom=2.5cm, left=2.5cm, right=2.5cm]{geometry}
\usepackage{xcolor}
\usepackage{listings}
\usepackage{tcolorbox}
\usepackage{tabularx}
\usepackage{booktabs}
\usepackage{enumitem}
\usepackage{titlesec}
\usepackage{fancyhdr}
\usepackage{hyperref}
\usepackage{tikz}
\usepackage{array}
\usepackage{amssymb}
\usepackage{amsmath}
\usepackage{multirow}
\usepackage{float}
\usepackage{setspace}
\usepackage{graphicx}
\usetikzlibrary{shapes.geometric, arrows.meta, positioning, calc, fit}

\definecolor{primary}{HTML}{1A237E}
\definecolor{accent}{HTML}{37474F}
\definecolor{success}{HTML}{2E7D32}
\definecolor{warning}{HTML}{E65100}
\definecolor{danger}{HTML}{B71C1C}
\definecolor{lightblue}{HTML}{E3F2FD}
\definecolor{lightgreen}{HTML}{E8F5E9}
\definecolor{lightorange}{HTML}{FFF3E0}
\definecolor{lightgray}{HTML}{F5F5F5}
\definecolor{darkgray}{HTML}{424242}
\definecolor{codebg}{HTML}{F8F8F8}
\definecolor{codeframe}{HTML}{CCCCCC}
\definecolor{nodecolor}{HTML}{BBDEFB}
\definecolor{topiccolor}{HTML}{C8E6C9}
\definecolor{hwcolor}{HTML}{FFE0B2}

\tcbuselibrary{skins, breakable}

\lstset{
  backgroundcolor=\color{codebg},
  basicstyle=\ttfamily\scriptsize,
  breaklines=true,
  frame=single,
  rulecolor=\color{codeframe},
  xleftmargin=4pt, xrightmargin=4pt, framexleftmargin=4pt,
  numbers=left, numberstyle=\tiny\color{darkgray},
  keywordstyle=\color{primary}\bfseries,
  commentstyle=\color{success}\itshape,
  stringstyle=\color{warning},
  showstringspaces=false,
}

\titleformat{\section}[block]
  {\color{primary}\Large\bfseries}{\thesection.}{8pt}{}
  [\vspace{2pt}\color{primary}\hrule\vspace{4pt}]

\titleformat{\subsection}[block]
  {\color{accent}\large\bfseries}{\thesubsection.}{6pt}{}

\titleformat{\subsubsection}[block]
  {\color{darkgray}\normalsize\bfseries}{\thesubsubsection.}{6pt}{}

\pagestyle{fancy}
\fancyhf{}
\fancyhead[L]{\textcolor{primary}{\textbf{Rapport de Stage} — Protova1 \& Protova2}}
\fancyhead[R]{\textcolor{darkgray}{\small SEBTI Douae}}
\fancyfoot[C]{\textcolor{darkgray}{\small \thepage}}
\renewcommand{\headrulewidth}{0.4pt}
\setstretch{1.15}

% ─────────────────────────────────────────────────────────────────
\begin{document}

% ======================== PAGE DE TITRE ========================
\begin{titlepage}
\thispagestyle{empty}
\begin{center}

\rule{\textwidth}{0.5pt}\\[0.4cm]
{\large\scshape Rapport de Stage de Fin d'Études}\\[0.2cm]
\rule{\textwidth}{0.5pt}

\vspace{2cm}

{\LARGE\bfseries Développement de Deux Plateformes de}\\[0.3cm]
{\LARGE\bfseries Véhicules RC Autonomes et Téléopérés}\\[0.6cm]
{\Large\bfseries Protova1 \textnormal{—} de zéro à la téléopération}\\[0.2cm]
{\Large\bfseries Protova2 \textnormal{—} démo de téléopération 5G}\\[0.4cm]
\rule{0.6\textwidth}{0.3pt}\\[0.4cm]
{\large Architecture, Implémentation, Problèmes et Perspectives}

\vspace{2cm}

\begin{tabular}{rl}
\textit{Auteure}              & SEBTI Douae \\[0.3cm]
\textit{Encadrant laboratoire} & Nicolas Baudesson \\[0.3cm]
\textit{Tuteur école}          & Samuel Dupont \\[0.3cm]
\textit{Durée du stage}        & Mars -- Juin 2026 (3 mois et 1 semaine) \\[0.3cm]
\textit{Plateformes}           & Jetson Nano, Raspberry Pi Pico RP2040 \\[0.3cm]
\textit{Environnement}         & ROS2 Foxy/Humble, FreeRTOS, Zenoh, 5G \\
\end{tabular}

\vspace{1.5cm}

\begin{minipage}{0.88\textwidth}
\small\itshape
Ce rapport présente les travaux réalisés durant un stage de fin d'études portant
simultanément sur deux plateformes de véhicules RC~: \textbf{Protova1}, reconstruite
intégralement depuis zéro avec intégration progressive des capteurs (encodeur, IMU,
LiDAR), et \textbf{Protova2}, plateforme de référence sur laquelle une démonstration de
téléopération complète via réseau 5G a été réalisée. Le rapport couvre l'architecture
matérielle et logicielle des deux systèmes, le détail des travaux chronologiques, les
problèmes rencontrés et leurs solutions, ainsi que les perspectives d'autonomie et
d'intégration 5G sur Protova1.
\end{minipage}

\vfill
\rule{\textwidth}{0.5pt}\\[0.3cm]
{\small Juin 2026}\\
\rule{\textwidth}{0.5pt}
\end{center}
\end{titlepage}

% ======================== RÉSUMÉ ========================
\newpage
\begin{tcolorbox}[colback=lightblue, colframe=primary, title={\textbf{Résumé}}]
Ce stage de 3 mois et 1 semaine a porté sur deux prototypes de véhicules RC autonomes
développés en parallèle au laboratoire. \textbf{Protova1} a été reconstruite de zéro~:
architecture Jetson Nano / Raspberry Pi Pico RP2040, pile ROS2 Foxy, FreeRTOS, téléopération
PS4 et G29, odométrie Ackermann, IMU BNO086, filtre anti-rebond hardware pour l'encodeur, et
arbitrage multi-sources via \texttt{twist\_mux}. \textbf{Protova2}, plateforme de référence déjà
opérationnelle, a fait l'objet d'une démonstration complète de téléopération via réseau 5G~:
intégration du dongle Quectel RM530N-GL, configuration réseau sur les deux extrémités (PC
opérateur et Jetson embarqué), middleware Zenoh pour traverser le réseau cellulaire, et
tableau de bord de monitoring en temps réel. Les performances mesurées atteignent 635~Mbps
en téléchargement et 70~Mbps en upload avec une latence inférieure à 30~ms. Lorsque
Protova1 sera complète, l'intégration 5G sera directement héritée de Protova2 — l'architecture
étant identique.
\end{tcolorbox}

\tableofcontents
\newpage

% =====================================================================
%                    1. INTRODUCTION
% =====================================================================
\section{Introduction}

\subsection{Contexte du projet}

Le projet \textbf{Protova} (Prototype de Véhicule Autonome) s'inscrit dans le cadre
de la recherche sur les véhicules autonomes à échelle réduite. Ces plateformes RC
constituent des outils pédagogiques et expérimentaux permettant de valider des
algorithmes de perception, de contrôle et de communication véhiculaire dans un
environnement maîtrisé.

Le stage s'est organisé autour de \textbf{deux plateformes distinctes travaillées
en parallèle}~:

\begin{itemize}
  \item \textbf{Protova1}~: nouveau prototype à construire intégralement depuis zéro,
        matériel neuf, logiciel à développer, capteurs à intégrer un par un.
  \item \textbf{Protova2}~: plateforme de référence déjà opérationnelle, sur laquelle
        la mission principale était de réaliser et valider la démonstration de
        téléopération complète via réseau 5G.
\end{itemize}

\subsection{Objectifs}

\begin{enumerate}[noitemsep]
  \item Reproduire sur Protova1 les fonctionnalités de Protova2~: téléopération
        multi-sources, odométrie, LiDAR, freinage automatique (AEB).
  \item Réaliser sur Protova2 la démo de téléopération 5G complète (dongle,
        réseau, Zenoh, monitoring).
  \item Préparer l'intégration 5G sur Protova1 — une fois celle-ci complète, le
        bridge Zenoh de Protova2 sera directement réutilisé.
\end{enumerate}

\subsection{Planning général}

\begin{center}
\begin{tabularx}{\textwidth}{@{}cXcc@{}}
\toprule
\textbf{Période} & \textbf{Travaux principaux} & \textbf{Prototype} & \textbf{Statut} \\
\midrule
Mois 1 & Architecture + Téléopération PS4 & Protova1 & \textcolor{success}{\textbf{Terminé}} \\
Mois 1 & Étude Protova2 + config réseau 5G & Protova2 & \textcolor{success}{\textbf{Terminé}} \\
Mois 2 & Odométrie + IMU + FreeRTOS & Protova1 & \textcolor{success}{\textbf{Terminé}} \\
Mois 2 & Zenoh + launch scripts + watchdog & Protova2 & \textcolor{success}{\textbf{Terminé}} \\
Mois 3 & G29 + twist\_mux + filtre encodeur & Protova1 & \textcolor{success}{\textbf{Terminé}} \\
Mois 3 & Démo 5G complète + monitoring & Protova2 & \textcolor{success}{\textbf{Terminé}} \\
Semaine+1 & LiDAR + AEB (Protova1) & Protova1 & \textcolor{warning}{\textbf{En cours}} \\
\bottomrule
\end{tabularx}
\end{center}


% =====================================================================
%                  2. PRÉSENTATION DES PLATEFORMES
% =====================================================================
\newpage
\section{Présentation des deux plateformes}

\subsection{Protova2 — Plateforme de référence}

Protova2 est un véhicule RC téléopéré opérationnel, équipé~:
\begin{itemize}[noitemsep]
  \item d'un \textbf{Jetson Nano} exécutant ROS2 Humble
  \item d'un \textbf{volant Logitech G29} avec retour de force (force feedback)
  \item d'un \textbf{LiDAR RPLidar} pour la détection d'obstacles
  \item d'une \textbf{caméra Orbbec Femto Bolt} (RGB-D)
  \item d'un \textbf{dongle 5G Quectel RM530N-GL} pour la communication longue portée
  \item d'un node \textbf{AEB} (Autonomous Emergency Braking) pour la sécurité
\end{itemize}

\subsection{Protova1 — Nouvelle plateforme construite de zéro}

Protova1 est construite sur le même modèle de châssis RC mais avec un matériel
entièrement neuf. Elle intègre~:
\begin{itemize}[noitemsep]
  \item un \textbf{Jetson Nano} (Ubuntu 22, ROS2 Foxy)
  \item un \textbf{Raspberry Pi Pico RP2040} (microcontrôleur FreeRTOS)
  \item un \textbf{encodeur KY-040} avec circuit de filtrage hardware
  \item une \textbf{centrale inertielle BNO086} (I2C)
  \item un \textbf{servo de direction} et un \textbf{ESC} (Electronic Speed Controller)
  \item une téléopération double-source~: \textbf{manette PS4} (Bluetooth) et
        \textbf{volant G29} (USB)
\end{itemize}

\begin{tcolorbox}[colback=lightgreen, colframe=success,
  title={\textbf{Lien entre les deux plateformes}}]
L'architecture logicielle de Protova1 est directement inspirée de Protova2.
Lorsque Protova1 sera complète (LiDAR + AEB), l'intégration 5G via Zenoh sera
\textbf{immédiatement réutilisable} depuis Protova2 sans développement supplémentaire —
les stacks ROS2, les topics et les launch scripts sont identiques.
\end{tcolorbox}


% =====================================================================
%              3. ARCHITECTURE PROTOVA1
% =====================================================================
\newpage
\section{Architecture de Protova1}

\subsection{Vue d'ensemble}

Protova1 repose sur une architecture à deux niveaux reliés par liaison série USB~:

\begin{center}
\begin{tikzpicture}[
  node distance=1.2cm and 2cm,
  box/.style={draw, rounded corners, minimum width=3cm, minimum height=0.85cm,
              font=\small, align=center},
  hw/.style={box, fill=hwcolor},
  sw/.style={box, fill=nodecolor},
  arr/.style={-{Stealth[length=2.5mm]}, thick},
]
\node[draw, dashed, fill=lightblue!20, minimum width=15cm, minimum height=7cm,
      rounded corners, label={[font=\bfseries\large]above:Jetson Nano — ROS2 Foxy}]
      (jetson) at (0,0) {};

\node[sw] (joy)  at (-5.5,  1.5) {joy\_node};
\node[sw] (ps4)  at (-2.5,  2.3) {ps4\_teleop};
\node[sw] (g29)  at (-2.5,  0.7) {G29Teleop};
\node[sw] (mux)  at ( 1.0,  1.5) {twist\_mux};
\node[sw] (tx)   at ( 4.5,  1.5) {TX};
\node[sw] (odom) at ( 1.0, -1.8) {ackermann\_odom};

\draw[arr] (joy) -- node[above,font=\tiny]{/joy} (ps4);
\draw[arr] (joy) -- node[below,font=\tiny]{/joy} (g29);
\draw[arr] (ps4) -- node[above,font=\tiny]{/cmd\_vel\_ps4} (mux);
\draw[arr] (g29) -- node[below,font=\tiny]{/cmd\_vel\_G29} (mux);
\draw[arr] (mux) -- node[above,font=\tiny]{/cmd\_vel} (tx);
\draw[arr, primary] (tx) |- node[right,font=\tiny,pos=0.3]{/velocity} (odom);

\node[hw] (pico) at (4.5, -2.5) {Pico RP2040\\FreeRTOS};
\draw[arr, warning, <->] (tx) -- node[right,font=\tiny]{Série USB 115200} (pico);
\end{tikzpicture}
\end{center}

\subsection{Niveau haut — Jetson Nano (ROS2 Foxy)}

\begin{center}
\begin{tabular}{llll}
\toprule
\textbf{Node} & \textbf{Fichier} & \textbf{Rôle} & \textbf{Statut} \\
\midrule
joy\_node & (package \texttt{joy}) & Lecture joystick USB/BT & \textcolor{success}{Terminé} \\
ps4\_teleop & ps4\_teleop.py & Mapping manette PS4 & \textcolor{success}{Terminé} \\
G29Teleop & G29Teleop.py & Mapping volant G29 + pédales & \textcolor{success}{Terminé} \\
twist\_mux & (ext.) & Arbitrage multi-sources & \textcolor{success}{Terminé} \\
TX & TX.py & Pont série bidirectionnel & \textcolor{success}{Terminé} \\
ackermann\_odom & ackermann\_odom.py & Odométrie Ackermann & \textcolor{success}{Terminé} \\
AEB & AEB.py & Freinage d'urgence automatique & \textcolor{warning}{Planifié} \\
dist\_min & dist\_min.py & Distance obstacle LiDAR & \textcolor{warning}{Planifié} \\
\bottomrule
\end{tabular}
\end{center}

\subsection{Niveau bas — Raspberry Pi Pico RP2040 (FreeRTOS)}

La migration vers FreeRTOS (depuis MicroPython) a été une décision technique majeure
pour garantir la déterminisme du contrôle temps réel.

\begin{center}
\begin{tabular}{lllll}
\toprule
\textbf{Tâche} & \textbf{Rôle} & \textbf{Priorité} & \textbf{Période} & \textbf{Trame} \\
\midrule
TaskSerialRX & Réception commandes & Haute & Événement & \texttt{CTRL,angle,throttle} \\
TaskControl  & PWM servo + ESC     & Haute & 20 ms & — \\
TaskEncoder  & Comptage ticks encodeur & Normale & 20 ms & \texttt{TICKS,N} \\
TaskIMU      & Lecture BNO086 & Normale & 50 ms & \texttt{IMU,qw,qx,qy,qz,yaw} \\
\bottomrule
\end{tabular}
\end{center}

\subsection{Brochage matériel}

\begin{center}
\begin{tabular}{llll}
\toprule
\textbf{GPIO} & \textbf{Périphérique} & \textbf{Signal} & \textbf{Protocole} \\
\midrule
GP2  & Encodeur KY-040 & CLK (via 74HC14) & Interruption FALLING \\
GP3  & Encodeur KY-040 & DT (via 74HC14)  & Interruption \\
GP4  & IMU BNO086      & SDA              & I2C0 (adresse 0x4B) \\
GP5  & IMU BNO086      & SCL              & I2C0 \\
GP15 & Servo direction & Signal PWM       & PWM 50 Hz \\
GP16 & ESC moteur      & Signal PWM       & PWM 50 Hz \\
\bottomrule
\end{tabular}
\end{center}

\subsection{Protocole de communication série}

\begin{center}
\begin{tabular}{lll}
\toprule
\textbf{Direction} & \textbf{Format} & \textbf{Exemple} \\
\midrule
Jetson → Pico & \texttt{CTRL,<angle>,<throttle>} & \texttt{CTRL,83,0.15} \\
Pico → Jetson & \texttt{TICKS,<delta>}            & \texttt{TICKS,3} \\
Pico → Jetson & \texttt{IMU,<qw>,<qx>,<qy>,<qz>,<yaw>} & \texttt{IMU,0.998,0.001,-0.003,0.052,6.0} \\
\bottomrule
\end{tabular}
\end{center}

\subsection{Modèle cinématique Ackermann}

Le node \texttt{ackermann\_odom} implémente le modèle bicyclette pour estimer la
position $(x, y, \theta)$ du véhicule. La vitesse de lacet est~:

\begin{equation}
  \omega = \frac{v \cdot \tan(\delta)}{L}
\end{equation}

avec $v$ la vitesse linéaire (issue de l'encodeur), $\delta$ l'angle de braquage et
$L = 0{,}245\ \text{m}$ l'empattement. L'intégration Euler à 50 Hz~:

\begin{align}
  \theta_{k+1} &= \theta_k + \omega \cdot \Delta t \\
  x_{k+1}      &= x_k + v \cdot \cos(\theta_k) \cdot \Delta t \\
  y_{k+1}      &= y_k + v \cdot \sin(\theta_k) \cdot \Delta t
\end{align}

La vitesse linéaire est calculée par le node TX à chaque trame encodeur~:

\begin{equation}
  v = \frac{N \times d_{\text{tick}}}{\Delta t}, \qquad
  d_{\text{tick}} = \frac{\pi \times D}{\text{ticks\_per\_rev} / \text{gear\_ratio}}
  = \frac{\pi \times 0{,}104}{49 / 0{,}902} = 6{,}01\ \text{mm/tick}
\end{equation}

\subsection{Arbitrage multi-sources — twist\_mux}

\begin{center}
\begin{tabular}{llll}
\toprule
\textbf{Source} & \textbf{Topic} & \textbf{Priorité} & \textbf{Timeout} \\
\midrule
AEB (sécurité) & /cmd\_vel\_aeb & 255 (maximale) & 5,0 s \\
PS4 Bluetooth  & /cmd\_vel\_ps4 & 60 & 0,5 s \\
G29 USB        & /cmd\_vel\_G29 & 30 & 0,5 s \\
\bottomrule
\end{tabular}
\end{center}

La priorité 255 de l'AEB garantit que le freinage d'urgence prévaut toujours sur
l'opérateur.

\subsection{Paramètres physiques calibrés}

\begin{center}
\begin{tabular}{llll}
\toprule
\textbf{Paramètre} & \textbf{Valeur} & \textbf{Méthode} \\
\midrule
Empattement               & 0,245 m  & Mesure directe \\
Diamètre roue             & 0,104 m  & Mesure directe \\
Ticks par tour            & 49       & Calibration Arduino \\
Rapport d'engrenage       & 0,902    & $15{,}7 / 17{,}4$ \\
Centre servo              & 83°      & Calibré sur sol plat \\
Amplitude servo           & 55°      & Mesure \\
Angle max braquage        & 0,5 rad (${\approx}28{,}6°$) & Mesure \\
\bottomrule
\end{tabular}
\end{center}


% =====================================================================
%              4. ARCHITECTURE PROTOVA2 — DÉMO 5G
% =====================================================================
\newpage
\section{Architecture de Protova2 — Démonstration 5G}

\subsection{Vue d'ensemble}

Protova2 est la plateforme de référence sur laquelle la démonstration de
téléopération 5G a été réalisée. Le système met en relation~:

\begin{center}
\begin{tikzpicture}[
  node distance=2cm and 3cm,
  box/.style={draw, rounded corners, minimum width=3.5cm, minimum height=1cm,
              font=\small, align=center},
  arr/.style={-{Stealth[length=3mm]}, very thick},
]
  \node[box, fill=nodecolor] (pc)    {PC Opérateur\\G29 + ROS2 Humble\\172.16.48.7};
  \node[box, fill=hwcolor,   right=5cm of pc] (car)
        {Jetson Nano (Voiture)\\ROS2 Humble\\172.16.48.6};
  \node[box, fill=lightblue, below=1.5cm of pc]  (dpc)  {Dongle 5G\\Quectel RM530N-GL};
  \node[box, fill=lightblue, below=1.5cm of car] (dcar) {Dongle 5G\\Quectel RM530N-GL};
  \node[box, fill=lightgreen, below=1.5cm of dpc, xshift=3.5cm]
        (zen) {Réseau 5G\\(Orange France)\\Zenoh bridge};

  \draw[arr] (pc)  -- (dpc);
  \draw[arr] (car) -- (dcar);
  \draw[arr, <->] (dpc) -- node[above, font=\small\bfseries]{5G} (dcar);
  \draw[arr, <->, dashed] (dpc)  -- (zen);
  \draw[arr, <->, dashed] (dcar) -- (zen);
\end{tikzpicture}
\end{center}

\subsection{Stack logiciel Protova2}

\begin{center}
\begin{tabular}{lll}
\toprule
\textbf{Composant} & \textbf{Rôle} & \textbf{Côté} \\
\midrule
ROS2 Humble + rmw\_zenoh\_cpp & Middleware DDS sur 5G & PC + Jetson \\
Zenoh router & Pont réseau 5G (port 7447) & PC (routeur) \\
G29 force feedback & Téléopération avec retour de force & PC \\
twist\_mux & Arbitrage PS4/G29/AEB & Jetson \\
TX.py / RX.py & Pont série Jetson ↔ actionneurs & Jetson \\
AEB + dist\_min & Freinage d'urgence LiDAR & Jetson \\
Orbbec Femto Bolt & Caméra RGB-D & Jetson \\
RPLidar & Détection obstacles & Jetson \\
monitor\_5g.py & Dashboard 5G temps réel (port 8088) & Jetson \\
bench\_5g\_live.py & Benchmark actif iperf3 (port 8089) & Jetson \\
\bottomrule
\end{tabular}
\end{center}

\subsection{Configuration réseau 5G}

Le dongle Quectel RM530N-GL (USB, VID \texttt{2c7c:0801}, réseau NR Band 78 TDD,
opérateur Bouygues Telecom France) est connecté aux deux machines.
Il apparaît sous Linux avec le préfixe \texttt{enx} (Ethernet-over-USB) ou renommé
\texttt{5g0} par une règle udev.

\begin{center}
\begin{tabular}{llll}
\toprule
\textbf{Machine} & \textbf{Interface} & \textbf{Adresse IP} & \textbf{Rôle Zenoh} \\
\midrule
PC opérateur     & \texttt{5g0}       & 172.16.48.7/28  & Hub (routeur central) \\
Jetson (voiture) & \texttt{5g0}       & 172.16.48.6/28  & Client (connect → PC) \\
Passerelle modem & —                  & 172.16.48.5     & Gateway APN opérateur \\
\bottomrule
\end{tabular}
\end{center}

\subsubsection{Renommage stable de l'interface — règle udev}

Par défaut, le dongle reçoit un nom basé sur son adresse MAC
(\texttt{enx66e352cb6c9c}) qui peut changer à chaque reconnexion USB.
Une règle udev déployée par \texttt{setup\_5g.sh} fixe un nom stable \texttt{5g0}~:

\begin{lstlisting}[language=bash]
# /etc/udev/rules.d/70-quectel-5g.rules
SUBSYSTEM=="net", ACTION=="add",
    ATTRS{idVendor}=="2c7c", ATTRS{idProduct}=="0801",
    NAME="5g0"
\end{lstlisting}

Le renommage prend effet après un redébranchement physique du dongle. Cette
configuration n'est réalisée qu'une seule fois.

\subsubsection{Bug /30 → /28 — adresse de broadcast}

Le modem Quectel distribue l'IP via DHCP avec un masque \textbf{/30}. Or dans un
bloc /30, les quatre adresses sont~:

\begin{center}
\begin{tabular}{ll}
\toprule
\textbf{Adresse} & \textbf{Rôle dans le /30} \\
\midrule
172.16.48.4 & Adresse réseau \\
172.16.48.5 & Gateway (modem) \\
172.16.48.6 & Jetson Nano (hôte) \\
172.16.48.7 & \textbf{Adresse broadcast} $\leftarrow$ or c'est l'IP fixe du PC ! \\
\bottomrule
\end{tabular}
\end{center}

Le noyau Linux refuse d'envoyer des paquets unicast vers l'adresse de broadcast.
Tout ping vers le PC retournait \textit{«~Do you want to ping broadcast?~»}.
La solution est d'élargir le masque à \textbf{/28} (broadcast = \texttt{.15})~:

\begin{lstlisting}[language=bash]
sudo ip addr del 172.16.48.6/30 brd 172.16.48.7 dev "$IFACE_5G"
sudo ip addr add 172.16.48.6/28 brd 172.16.48.15 dev "$IFACE_5G"
\end{lstlisting}

Cette correction est intégrée dans \texttt{launch\_car.sh} et s'applique
automatiquement à chaque démarrage.

\subsection{Zenoh — bridge ROS2 sur réseau cellulaire}

Zenoh remplace le DDS natif de ROS2 car ce dernier repose sur le multicast UDP,
incompatible avec les réseaux cellulaires (pas de multicast sur la 5G).

\subsubsection{Architecture hub-and-spoke}

L'architecture adoptée est de type \textbf{étoile (hub-and-spoke) sur TCP}~:

\begin{itemize}[noitemsep]
  \item \textbf{Routeur PC} (\texttt{zenoh\_router\_pc.json5})~: écoute sur
        \texttt{0.0.0.0:7447}, n'initie aucune connexion sortante.
  \item \textbf{Routeur Jetson} (\texttt{zenoh\_router\_car.json5})~: écoute sur
        \texttt{0.0.0.0:7447} et se connecte vers \texttt{tcp/172.16.48.7:7447}.
        C'est le Jetson qui \textbf{ouvre la connexion TCP} à travers le lien 5G.
  \item \textbf{Nodes ROS2}~: se connectent en mode \emph{client} à leur routeur
        local (\texttt{127.0.0.1:7447}) via \texttt{zenoh\_session\_car.json5}.
        Aucun node ne communique directement en 5G.
\end{itemize}

Configurations JSON5 utilisées~:

\begin{lstlisting}[language=bash]
// zenoh_router_pc.json5 -- hub central (PC)
{ mode: "router",
  listen:  { endpoints: ["tcp/0.0.0.0:7447"] },
  scouting: { multicast: { enabled: false } } }

// zenoh_router_car.json5 -- Jetson Nano
{ mode: "router",
  listen:  { endpoints: ["tcp/0.0.0.0:7447"] },
  connect: { endpoints: ["tcp/172.16.48.7:7447"] },
  scouting: { multicast: { enabled: false } } }

// zenoh_session_car.json5 -- sessions ROS2 locales
{ mode: "client",
  connect: { endpoints: ["tcp/127.0.0.1:7447"] },
  scouting: { multicast: { enabled: false } } }
\end{lstlisting}

\textbf{Vérification de l'appairage} (depuis le Jetson)~:
\begin{lstlisting}[language=bash]
ss -tn | grep 172.16.48.7:7447   # doit afficher ESTAB
\end{lstlisting}

\begin{center}
\begin{tabular}{lll}
\toprule
\textbf{Variable d'environnement} & \textbf{Valeur} & \textbf{Rôle} \\
\midrule
\texttt{RMW\_IMPLEMENTATION} & \texttt{rmw\_zenoh\_cpp} & Active Zenoh \\
\texttt{ROS\_DOMAIN\_ID}     & 125 & Isolation du domaine ROS \\
\texttt{ZENOH\_SESSION\_CONFIG\_URI} & \texttt{\$HOME/zenoh\_session\_car.json5} & Config session \\
\texttt{ZENOH\_ROUTER\_CONFIG\_URI}  & \texttt{\$HOME/zenoh\_router\_car.json5} & Config routeur \\
\bottomrule
\end{tabular}
\end{center}

\subsection{Script de démarrage — launch\_car.sh}

Le script \texttt{launch\_car.sh} orchestre le démarrage complet du côté voiture~:
interface 5G, IP statique, watchdog réseau, routeur Zenoh, tous les nodes ROS2, et
le dashboard de monitoring. Les étapes principales~:

\begin{enumerate}[noitemsep]
  \item Détection automatique du dongle 5G (udev ou VID Quectel \texttt{2c7c})
  \item Attribution de l'IP statique \texttt{172.16.48.6/28}
  \item Lancement du watchdog (réapplique la config si le dongle rebondit sur USB)
  \item Démarrage du routeur Zenoh (\texttt{rmw\_zenohd})
  \item Lancement des nodes~: \texttt{g29\_teleop}, \texttt{TX}, \texttt{RX},
        \texttt{twist\_mux}, caméra Orbbec
  \item Démarrage de \texttt{monitor\_5g.py} (dashboard web \texttt{:8088})
\end{enumerate}

\subsection{Procédure de diagnostic rapide}

En cas de panne du lien, l'ordre de vérification recommandé~:

\begin{enumerate}[noitemsep]
  \item \texttt{ip -br addr show 5g0} — l'interface a-t-elle bien
        \texttt{172.16.48.6/28} (et non /30) ?
  \item \texttt{ping 172.16.48.7} — le lien IP est-il actif ?
  \item \texttt{ss -tn | grep 172.16.48.7:7447} — les routeurs Zenoh sont-ils
        appairés ? (panne la plus fréquente)
  \item \texttt{ros2 topic info /cmd\_vel\_G29 -v} — nombre de publishers $> 0$ ?
\end{enumerate}

\subsection{Outils de monitoring développés}

Deux scripts Python ont été développés pour monitorer le lien 5G~:

\subsubsection{monitor\_5g.py — Monitoring passif temps réel}

Tourne en permanence sur le Jetson, \textbf{sans générer de trafic}. Lit les compteurs
sysfs du dongle (\texttt{/sys/class/net/<iface>/statistics/}) toutes les 0,5~s et
sert un dashboard web sur le port 8088.

\textbf{Métriques collectées}~: latence ICMP, RTT TCP Zenoh (via \texttt{ss}),
débit RX/TX passif (Mbps), paquets/s, perte paquets, fréquence topics ROS,
CPU Jetson, état Zenoh.

\textbf{Formule du débit passif}~:
\begin{equation}
  D_{Mbps} = \frac{(B_k - B_{k-1}) \times 8}{\Delta t \times 10^6}
\end{equation}

\subsubsection{bench\_5g\_live.py — Benchmark actif iperf3}

Effectue des tests \texttt{iperf3} périodiques (toutes les 30~s, 5~s par direction)
pour mesurer la \textbf{capacité réelle} du lien 5G. Dashboard sur le port 8089 avec
un bouton «~Tester maintenant~».

\textbf{Métriques}~: débit TCP UL (Jetson→PC) et DL (PC→Jetson) en Mbps,
latence ICMP continue, trafic passif.

\begin{equation}
  D_{iperf,Mbps} = \frac{\texttt{bits\_per\_second}}{10^6}
  = \frac{8 \times \texttt{bytes\_transferred}}{D_{test} \times 10^6}
\end{equation}

\subsection{Performances mesurées}

\begin{center}
\begin{tabular}{lll}
\toprule
\textbf{Métrique} & \textbf{Valeur mesurée} & \textbf{Outil} \\
\midrule
Débit téléchargement (5G) & 635,52 Mbps & Speedtest.net \\
Débit upload (5G)         & 70,25 Mbps  & Speedtest.net \\
Latence ICMP (RTT)        & $\approx$ 23--30 ms & monitor\_5g.py \\
Perte paquets             & $< 1\%$ & monitor\_5g.py \\
Opérateur                 & Orange France (Lyon) & — \\
\bottomrule
\end{tabular}
\end{center}


% =====================================================================
%         5. PROBLÈMES RENCONTRÉS ET SOLUTIONS
% =====================================================================
\newpage
\section{Problèmes rencontrés et solutions}

Cette section recense les principaux obstacles techniques rencontrés et les solutions
apportées, sur les deux plateformes.

\subsection{Problèmes sur Protova1}

\subsubsection{Migration MicroPython → FreeRTOS}

\textbf{Problème~:} MicroPython sur le Pico ne permettait pas de gérer simultanément
les interruptions d'encodeur (très haute fréquence), la lecture I2C de l'IMU et la
communication série sans pertes.

\textbf{Solution~:} Migration complète vers Arduino + FreeRTOS avec 4 tâches concurrentes
et priorités définies. Le déterminisme temporel est garanti par les mécanismes de
scheduling de FreeRTOS.

\subsubsection{Rebonds de l'encodeur KY-040}

\textbf{Problème~:} L'encodeur KY-040 produit des signaux parasites (rebonds) lors de
chaque transition mécanique, causant un comptage erroné des ticks et une odométrie
fausse.

\textbf{Solution~:} Conception d'un circuit de filtrage hardware en deux étages~:
\begin{enumerate}[noitemsep]
  \item Filtre RC passe-bas ($\tau = R \times C = 0{,}1\ \text{ms}$) qui atténue les
        transitoires haute fréquence.
  \item Trigger de Schmitt 74HC14 qui restitue un signal carré propre grâce à
        l'hystérésis de ses deux seuils de commutation.
\end{enumerate}
Le 74HC14 inversant le signal, l'interruption est configurée en \texttt{FALLING}
au lieu de \texttt{RISING}.

\subsubsection{Adresse I2C non standard de l'IMU BNO086}

\textbf{Problème~:} Le BNO086 répond à l'adresse \texttt{0x4B} (broche ADR à +3V3)
et non à l'adresse par défaut \texttt{0x28}, ce qui n'était pas documenté dans les
exemples disponibles.

\textbf{Solution~:} Identification de l'adresse par scan I2C sur le Pico
(\texttt{i2c.scan()}), puis configuration explicite dans la bibliothèque BNO08x.

\subsubsection{Saturation du volant G29}

\textbf{Problème~:} Le G29 tourne sur 2,5 tours complets, mais la voiture RC ne
nécessite qu'un quart de tour. Les 75\% restants étaient inutilisables.

\textbf{Solution~:} Saturation logicielle dans \texttt{G29Teleop.py}~:
\begin{equation}
  \text{steer\_scaled} = \text{clamp}\!\left(\frac{\text{steer\_raw}}{0{,}33},\ -1,\ +1\right)
\end{equation}
À partir de 33\% de la course, les roues sont en butée — comportement identique à
celui de Protova2.

\subsection{Problèmes sur Protova2 — Réseau 5G}

\subsubsection{Masque de sous-réseau /30 — PC injoignable}

\textbf{Problème~:} Le modem Quectel, en mode ECM/NCM, distribue l'adresse IP
via DHCP avec un masque \texttt{/30}. Dans un bloc /30, l'adresse
\texttt{172.16.48.7} est l'adresse de \textbf{broadcast} — or c'est précisément
l'IP fixe du PC opérateur. Tout ping ou connexion vers le PC retournait~:
\begin{lstlisting}[language=bash]
Do you want to ping broadcast? Then -b. Check your local firewall rules.
\end{lstlisting}

\textbf{Cause racine~:} L'APN opérateur utilise réellement un sous-réseau \texttt{/28}
partagé entre les deux SIM. Le modem annonce cependant un \texttt{/30} au niveau
du lien USB Ethernet, ce qui masque la portée réelle du réseau.

\textbf{Solution~:} Réécrire le préfixe IP de \texttt{/30} à \texttt{/28} dans
\texttt{launch\_car.sh}~:
\begin{lstlisting}[language=bash]
sudo ip addr del 172.16.48.6/30 brd 172.16.48.7 dev "$IFACE_5G" 2>/dev/null
sudo ip addr add 172.16.48.6/28 brd 172.16.48.15 dev "$IFACE_5G"
\end{lstlisting}
Avec \texttt{/28}, l'adresse broadcast est \texttt{172.16.48.15}~: les adresses
\texttt{.6} et \texttt{.7} deviennent deux hôtes unicast valides sur le même réseau.

\subsubsection{Diagnostic modem par commandes AT}

Lors de l'initialisation, le dongle ne répondait pas au DHCP après la mise sous
tension. Un diagnostic via commandes AT (Python, \texttt{serial.Serial} 115200 bauds
sur \texttt{/dev/ttyUSB3}) a permis d'identifier l'état du modem~:

\begin{center}
\begin{tabular}{lll}
\toprule
\textbf{Commande AT} & \textbf{Réponse} & \textbf{Interprétation} \\
\midrule
\texttt{AT+COPS?}   & Bouygues Telecom, 11 & Réseau NR (5G), enregistré \\
\texttt{AT+CGPADDR} & 172.16.48.6          & Contexte PDP actif, IP assignée \\
\texttt{AT+CSQ}     & 99,99                & Signal non disponible (ECM actif) \\
\bottomrule
\end{tabular}
\end{center}

Le modem avait bien une IP mais le pont USB Ethernet n'avait pas démarré son
serveur DHCP. La réinitialisation douce a résolu le problème~:
\begin{lstlisting}[language=python]
s.write(b'AT+CFUN=1,1\r\n')   # reset logiciel du modem
time.sleep(8)                   # attente redemarrage
\end{lstlisting}
Après reset, le noyau a détecté une nouvelle interface USB Ethernet et le DHCP
a répondu normalement.

\subsubsection{Route par défaut DHCP — Perte de connexion internet}

\textbf{Problème~:} Le dongle 5G pousse via DHCP une route par défaut via
\texttt{172.16.48.8} avec une métrique de 100, inférieure à celle du WiFi (600).
Résultat~: tout le trafic internet passait par la 5G, ralentissant les deux.

Le script \texttt{launch\_car.sh} tentait de supprimer cette route avec~:
\begin{lstlisting}[language=bash]
sudo ip route del default via 172.16.48.5   # ne fonctionnait pas : mauvaise IP
\end{lstlisting}
La commande échouait silencieusement car la gateway réelle était \texttt{172.16.48.8}.

\textbf{Solution~:} Suppression par interface plutôt que par gateway~:
\begin{lstlisting}[language=bash]
sudo ip route del default dev "$IFACE_5G"   # fonctionne quelle que soit la gateway
\end{lstlisting}

\subsubsection{Dongle 5G qui se déconnecte de l'USB}

\textbf{Problème~:} Le dongle Quectel rebondit occasionnellement sur le bus USB,
perdant son adresse IP et coupant la liaison 5G sans préavis.

\textbf{Solution~:} Mise en place d'un \textbf{watchdog réseau} dans
\texttt{launch\_car.sh}, tournaant en fond en boucle toutes les 5~s~:
\begin{lstlisting}[language=bash]
while true; do
    if ip link show 5g0 &>/dev/null && \
       ! ip -4 addr show 5g0 | grep -q "172.16.48.6/28"; then
        sudo ip link set 5g0 up
        sudo ip addr add 172.16.48.6/28 brd 172.16.48.15 dev 5g0
        echo "[watchdog] IP config reapplied"
    fi
    sleep 5
done
\end{lstlisting}

\subsubsection{Variables d'environnement Zenoh incorrectes}

\textbf{Problème~:} La documentation Zenoh indique \texttt{ZENOH\_CONFIG} mais
\texttt{rmw\_zenoh\_cpp} (version ROS2 Humble) utilise en réalité
\texttt{ZENOH\_SESSION\_CONFIG\_URI} et \texttt{ZENOH\_ROUTER\_CONFIG\_URI}.
Les nodes démarraient sans config Zenoh et ne se connectaient pas.

\textbf{Solution~:} Identification des bonnes variables dans le code source de
\texttt{rmw\_zenoh\_cpp} et mise à jour de tous les scripts de lancement.

\subsubsection{Fichiers de configuration Zenoh en doublon}

\textbf{Problème~:} Des fichiers de configuration Zenoh existaient à plusieurs
emplacements (\texttt{\textasciitilde/protova2/} et \texttt{\textasciitilde/}), avec
des contenus contradictoires. Le fichier
\texttt{\textasciitilde/protova2/zenoh\_session\_pc.json5} contenait en réalité
une configuration de \textit{routeur} et non de \textit{session}.

\textbf{Solution~:} Identification et documentation des fichiers actifs~:
seuls \texttt{\textasciitilde/zenoh\_session\_car.json5} et
\texttt{\textasciitilde/zenoh\_session\_pc.json5} (répertoire home) sont utilisés
par les scripts. Les fichiers dans \texttt{\textasciitilde/protova2/} sont des
artéfacts non utilisés.

\subsubsection{Unités de mesure incorrectes dans monitor\_5g.py (kb/s vs Mbps)}

\textbf{Problème 1~:} Le dashboard affichait le débit en kbit/s alors que le contexte
(lien 5G, référence speedtest à 635 Mbps) nécessitait des Mbit/s.

\textbf{Problème 2 (découlant de la correction)~:} En changeant l'unité de kb/s à
Mbps, la formule Python devenait~:
\begin{lstlisting}[language=python]
rx = max(0, cur[0] - prev[0]) * 8 / dt / 1_000_000  # Mbps
\end{lstlisting}
Mais le stockage des échantillons utilisait \texttt{rnd(rx, 1)}, soit un arrondi à
1 décimale. Le trafic normal (~0,022 Mbps) était arrondi à \texttt{0.0} — le
dashboard affichait 0 en permanence.

\textbf{Solution~:} Augmentation de la précision à 4 décimales~:
\begin{lstlisting}[language=python]
'rx': rnd(rx, 4), 'tx': rnd(tx, 4),   # 0.0220 Mbps = 22 kb/s resolu
\end{lstlisting}
Et ajout d'un auto-scaling JavaScript~: valeurs $< 1$ Mbps affichées en kb/s,
valeurs $\geq 1$ Mbps en Mbps.

\subsubsection{Piège SSH — coupure WiFi tue toute la pile}

\textbf{Problème~:} Lors de la démo sans WiFi, \texttt{launch\_car.sh} tournait
dans une session SSH ouverte via le WiFi. En coupant le WiFi pour basculer
complètement sur la 5G, la session SSH se fermait brusquement — et le
mécanisme \texttt{trap + wait} du script envoyait un signal de terminaison à tous
les processus enfants (nodes ROS2, routeur Zenoh, watchdog). La pile entière
s'arrêtait au moment précis où la démo devait commencer.

\textbf{Solution~:} Relancer la pile depuis une \textbf{session SSH ouverte via
l'IP 5G} (\texttt{ssh protova2@172.16.48.6}), en mode détaché de la session~:
\begin{lstlisting}[language=bash]
nohup ~/launch_car.sh > ~/launch_car.log 2>&1 & disown
\end{lstlisting}
Le processus \texttt{disown} détache le script du terminal SSH~: même si la
connexion se ferme, la pile continue de tourner.

\textbf{Procédure retenue pour la démo sans WiFi}~:
\begin{enumerate}[noitemsep]
  \item Démarrer les deux dongles + vérifier le ping 5G (WiFi encore actif)
  \item Ouvrir une session SSH via l'IP 5G~: \texttt{ssh protova2@172.16.48.6}
  \item Relancer la pile en mode détaché (commande ci-dessus)
  \item Couper le WiFi sur le Jetson~: \texttt{sudo nmcli radio wifi off}
  \item Couper le WiFi sur le PC~: \texttt{nmcli radio wifi off}
  \item Vérifier~: \texttt{ss -tn | grep 7447} doit afficher \texttt{ESTAB}
\end{enumerate}

\subsubsection{Tests iperf3 silencieusement échoués}

\textbf{Problème~:} Les tests iperf3 périodiques dans le dashboard retournaient
systématiquement \texttt{None} sans message d'erreur visible. Le bouton
«~Tester~» semblait ne rien faire.

\textbf{Cause~:} \texttt{iperf3 -s} n'était pas lancé sur le PC opérateur. Le client
iperf3 sur le Jetson échouait à se connecter, mais l'erreur était capturée
silencieusement dans \texttt{last\_err} sans être affichée.

\textbf{Solution~:} Affichage explicite des erreurs dans le dashboard
(\texttt{iperf\_err} en rouge), documentation du prérequis~:
\texttt{iperf3 -s} doit tourner sur le PC avant tout test.


% =====================================================================
%                    6. TRAVAUX RÉALISÉS
% =====================================================================
\newpage
\section{Bilan des travaux réalisés}

\subsection{Protova1 — État d'avancement}

\begin{center}
\begin{tabularx}{\textwidth}{@{}cXcc@{}}
\toprule
\textbf{Étape} & \textbf{Description} & \textbf{Statut} & \textbf{\%} \\
\midrule
1 & Migration MicroPython → FreeRTOS & \textcolor{success}{\textbf{Terminé}} & 100 \\
2 & Architecture Jetson + Pico & \textcolor{success}{\textbf{Terminé}} & 100 \\
3 & Téléopération PS4 Bluetooth & \textcolor{success}{\textbf{Terminé}} & 100 \\
4 & Odométrie Ackermann + encodeur KY-040 & \textcolor{success}{\textbf{Terminé}} & 100 \\
5 & IMU BNO086 (I2C + quaternions) & \textcolor{success}{\textbf{Terminé}} & 100 \\
6 & Circuit filtre encodeur (RC + 74HC14) & \textcolor{success}{\textbf{Terminé}} & 100 \\
7 & Téléopération G29 + pédales & \textcolor{success}{\textbf{Terminé}} & 100 \\
8 & Arbitrage twist\_mux multi-sources & \textcolor{success}{\textbf{Terminé}} & 100 \\
9 & LiDAR + détection obstacles & \textcolor{warning}{\textbf{En cours}} & 30 \\
10 & AEB (freinage d'urgence) & \textcolor{warning}{\textbf{Planifié}} & 0 \\
11 & 5G via Zenoh & \textit{Hérité de Protova2} & — \\
\midrule
\multicolumn{2}{l}{\textbf{Avancement global Protova1}} & \multicolumn{2}{c}{\textbf{85\%}} \\
\bottomrule
\end{tabularx}
\end{center}

\subsection{Protova2 — Démo 5G}

\begin{center}
\begin{tabularx}{\textwidth}{@{}cXc@{}}
\toprule
\textbf{Étape} & \textbf{Description} & \textbf{Statut} \\
\midrule
1 & Étude architecture existante Protova2 & \textcolor{success}{\textbf{Terminé}} \\
2 & Configuration dongle 5G (IP, masque /28) & \textcolor{success}{\textbf{Terminé}} \\
3 & Configuration Zenoh (routeur PC, session Jetson) & \textcolor{success}{\textbf{Terminé}} \\
4 & Script launch\_car.sh complet + watchdog & \textcolor{success}{\textbf{Terminé}} \\
5 & Démo téléopération G29 via 5G fonctionnelle & \textcolor{success}{\textbf{Terminé}} \\
6 & monitor\_5g.py (dashboard passif temps réel) & \textcolor{success}{\textbf{Terminé}} \\
7 & bench\_5g\_live.py (benchmark actif iperf3) & \textcolor{success}{\textbf{Terminé}} \\
8 & Validation performances 5G (635 Mbps / 30 ms) & \textcolor{success}{\textbf{Terminé}} \\
\midrule
\multicolumn{2}{l}{\textbf{Avancement global Protova2 5G}} & \textbf{100\%} \\
\bottomrule
\end{tabularx}
\end{center}


% =====================================================================
%                    7. PERSPECTIVES
% =====================================================================
\newpage
\section{Perspectives}

\subsection{Complétion de Protova1}

Les prochaines étapes à court terme pour finaliser Protova1 sont~:

\begin{enumerate}
  \item \textbf{Intégration LiDAR RPLidar}~: lecture des données de scan, publication sur
        \texttt{/scan} (sensor\_msgs/LaserScan).
  \item \textbf{Node dist\_min}~: extraction de la distance minimale dans un cône frontal
        de $\pm 10°$.
  \item \textbf{Node AEB}~: calcul du TTC (Time To Collision) et injection d'une commande
        d'arrêt sur \texttt{/cmd\_vel\_aeb} (priorité 255 dans twist\_mux)~:
        \begin{equation}
          \text{TTC} = \frac{d_{\text{front}}}{v}, \quad \text{si TTC} < 1{,}8\ \text{s}
          \Rightarrow \text{freinage d'urgence}
        \end{equation}
  \item \textbf{Fusion IMU dans l'odométrie}~: remplacement du $\theta$ géométrique
        par le yaw direct de l'IMU BNO086~:
        \begin{equation}
          \theta_{IMU} = \text{atan2}\!\left(2(q_w q_z + q_x q_y),\;
          1 - 2(q_y^2 + q_z^2)\right)
        \end{equation}
\end{enumerate}

\subsection{Intégration 5G sur Protova1}

Une fois Protova1 complète, la communication 5G sera directement héritée de Protova2.
Les éléments à transférer sont~:
\begin{itemize}[noitemsep]
  \item Le script \texttt{launch\_car.sh} (avec watchdog et config IP)
  \item Les fichiers de configuration Zenoh (\texttt{zenoh\_session\_car.json5})
  \item Les outils de monitoring (\texttt{monitor\_5g.py}, \texttt{bench\_5g\_live.py})
\end{itemize}

Aucun développement logiciel supplémentaire n'est nécessaire~: l'architecture
ROS2 de Protova1 est identique à celle de Protova2.

\subsection{Vers l'autonomie}

À plus long terme, trois niveaux d'autonomie sont envisagés~:

\begin{center}
\begin{tabularx}{\textwidth}{@{}cXcl@{}}
\toprule
\textbf{Niveau} & \textbf{Description} & \textbf{Technologies} & \textbf{Complexité} \\
\midrule
1 & Suivi de couloir LiDAR + PID & LiDAR + PID & Faible \\
2 & Circuit défini par waypoints & Pure Pursuit + AEB & Moyenne \\
3 & Navigation SLAM sans GPS & slam\_toolbox + Nav2 + EKF & Élevée \\
\bottomrule
\end{tabularx}
\end{center}

L'algorithme Pure Pursuit (niveau 2) calcule l'angle de braquage vers un point
cible (lookahead point) à distance $L_d$~:
\begin{equation}
  \delta = \arctan\!\left(\frac{2 L \sin(\alpha)}{L_d}\right)
\end{equation}

où $\alpha$ est l'angle entre le cap du véhicule et le waypoint cible, et $L$
l'empattement.


% =====================================================================
%                    8. CONCLUSION
% =====================================================================
\newpage
\section{Conclusion}

Ce stage de trois mois et une semaine a permis de mener de front deux projets
complémentaires et exigeants sur le plan technique.

Sur \textbf{Protova1}, les fondations techniques sont solidement établies~: architecture
à deux niveaux (Jetson Nano / Raspberry Pi Pico), pile ROS2 Foxy complète, four
tâches FreeRTOS, téléopération double-source (PS4 et G29), arbitrage multi-sources
avec sécurité intégrée, odométrie Ackermann à 50~Hz, IMU BNO086, et circuit de
filtrage hardware pour l'encodeur. La téléopération G29 et le twist\_mux sont
pleinement opérationnels.

Sur \textbf{Protova2}, la démonstration de téléopération via réseau 5G est complète
et validée~: le lien Jetson ↔ PC opérateur fonctionne via les dongles Quectel
RM530N-GL avec une latence inférieure à 30~ms et un débit disponible de 635~Mbps.
Le middleware Zenoh permet à ROS2 de traverser le réseau cellulaire de manière
transparente pour les applications. Deux outils de monitoring dédiés ont été
développés~: \texttt{monitor\_5g.py} pour l'observation passive en temps réel, et
\texttt{bench\_5g\_live.py} pour la mesure active de la capacité du lien.

La modularité des architectures choisies — nodes ROS2 indépendants, topics
standardisés, scripts de lancement paramétrables — garantit une réutilisation
directe entre les deux plateformes. L'intégration 5G sur Protova1, dès que celle-ci
sera complète, ne nécessitera aucun développement supplémentaire.

\vspace{1cm}

\begin{tcolorbox}[colback=lightblue, colframe=primary]
\textbf{Résumé des contributions principales}
\begin{itemize}[noitemsep]
  \item Architecture complète et fonctionnelle de Protova1 (hardware + software)
  \item Démonstration 5G opérationnelle sur Protova2
  \item Deux outils de monitoring 5G (passif et actif)
  \item Résolution de 7 problèmes techniques documentés
  \item Base technique prête pour l'autonomie (LiDAR + AEB + navigation SLAM)
\end{itemize}
\end{tcolorbox}

\end{document}
